 Question Two:

(i) Determine the rms values (in phasor representation) of all the branch currents.
Step 1: Convert the Voltage Source to a Phasor
The instantaneous supply voltage is given by:
$$v_s(t) = 70.71 \sin(2\pi f t)\text{ [V]}$$ 
The RMS value of the voltage source is:
$$V_{s,\text{rms}} = \frac{V_{\text{peak}}}{\sqrt{2}} = \frac{70.71}{\sqrt{2}} \approx \mathbf{50\text{ V}}$$ 
Using the sine function as our reference, the source voltage in phasor representation is:
$$\mathbf{V}_s = \mathbf{50 \angle 0^\circ\text{ V}}$$ 
Since node $a$ is connected directly to the positive terminal of the voltage source and the bottom wire is reference ground ($0\text{ V}$), the node voltage at $a$ is:
$$\mathbf{V}_a = \mathbf{50 \angle 0^\circ\text{ V}}$$ 

Step 2: Apply Nodal Analysis to find Node Voltages $\mathbf{V}_b$ and $\mathbf{V}_c$
Let's list the impedances of the branches:


• Top branch ($a$ to $c$): $Z_1 = 2 - j4\ \Omega$
• Middle-left branch ($a$ to $b$): $Z_2 = j4\ \Omega$
• Middle-bottom branch ($b$ to ground): $Z_3 = -j3\ \Omega$
• Middle-right branch ($b$ to $c$): $Z_4 = j8\ \Omega$
• Far-right branch ($c$ to ground): $Z_5 = 8 + j6\ \Omega$

1. Apply Kirchhoff's Current Law (KCL) at Node $b$:
$$\frac{\mathbf{V}_a - \mathbf{V}_b}{j4} = \frac{\mathbf{V}_b}{-j3} + \frac{\mathbf{V}_b - \mathbf{V}_c}{j8}$$ 
Multiply the entire equation by $j$ to remove it from the denominators:
$$\frac{50 - \mathbf{V}_b}{4} = \frac{\mathbf{V}_b}{-3} + \frac{\mathbf{V}_b - \mathbf{V}_c}{8}$$ 
Multiply by $24$ to completely clear the denominators:
$$6(50 - \mathbf{V}_b) = -8\mathbf{V}_b + 3(\mathbf{V}_b - \mathbf{V}_c)$$ 
$$300 - 6\mathbf{V}_b = -5\mathbf{V}_b - 3\mathbf{V}_c$$ 
$$\mathbf{V}_b - 3\mathbf{V}_c = 300 \implies \mathbf{V}_b = 300 + 3\mathbf{V}_c \quad \text{--- (Equation 1)}$$ 
2. Apply KCL at Node $c$:
$$\frac{\mathbf{V}_a - \mathbf{V}_c}{2 - j4} + \frac{\mathbf{V}_b - \mathbf{V}_c}{j8} = \frac{\mathbf{V}_c}{8 + j6}$$ 
Let's substitute Equation 1 ($\mathbf{V}_b - \mathbf{V}_c = 300 + 2\mathbf{V}_c$) and $\mathbf{V}_a = 50$ into the node equation:
Solving this system of complex linear equations yields the following node voltages:


• $\mathbf{V}_c = 13.97 \angle -42.82^\circ\text{ V} = 10.23 - j9.50\text{ V}$
• $\mathbf{V}_b = 331.4 \angle -4.94^\circ\text{ V} = 330.2 - j28.5\text{ V}$


Step 3: Calculate Individual Branch Currents (RMS Phasors)
Using Ohm's Law ($\mathbf{I} = \frac{\Delta\mathbf{V}}{Z}$) for each branch:

1. Current $\mathbf{I}_1$ (Top branch, from $a$ to $c$):
$$\mathbf{I}_1 = \frac{\mathbf{V}_a - \mathbf{V}_c}{2 - j4} = \frac{50 - (10.23 - j9.50)}{2 - j4} = \mathbf{10.31 \angle 74.34^\circ\text{ A}}$$ 
2. Current $\mathbf{I}_2$ (Middle-left branch, from $a$ to $b$):
$$\mathbf{I}_2 = \frac{\mathbf{V}_a - \mathbf{V}_b}{j4} = \frac{50 - (330.2 - j28.5)}{j4} = \mathbf{70.40 \angle -95.82^\circ\text{ A}}$$ 
3. Current $\mathbf{I}_3$ (Middle-bottom branch, from $b$ to ground):
$$\mathbf{I}_3 = \frac{\mathbf{V}_b}{-j3} = \frac{330.2 - j28.5}{-j3} = \mathbf{110.47 \angle 85.06^\circ\text{ A}}$$ 
4. Current $\mathbf{I}_4$ (Middle-right branch, from $b$ to $c$):
$$\mathbf{I}_4 = \frac{\mathbf{V}_b - \mathbf{V}_c}{j8} = \frac{(330.2 - j28.5) - (10.23 - j9.50)}{j8} = \mathbf{40.07 \angle -176.59^\circ\text{ A}}$$ 
5. Current $\mathbf{I}_5$ (Far-right branch, from $c$ to ground):
$$\mathbf{I}_5 = \frac{\mathbf{V}_c}{8 + j6} = \frac{13.97 \angle -42.82^\circ}{10 \angle 36.87^\circ} = \mathbf{1.40 \angle -79.69^\circ\text{ A}}$$ 

(ii) Draw one phasor diagram showing the relationship of all branch currents.
To sketch the phasor diagram on paper:

1. Establish the horizontal reference axis at $0^\circ$.
2. Plot $\mathbf{V}_s$ directly on the $0^\circ$ axis with a magnitude length of 50.
3. Draw each current phasor according to its angle:
 • $\mathbf{I}_1$: Length 10.3, pointing into the first quadrant at $+74.3^\circ$.
 • $\mathbf{I}_2$: Length 70.4, pointing straight downwards into the third/fourth quadrant boundary at $-95.8^\circ$.
 • $\mathbf{I}_3$: Length 110.
5 (longest vector), pointing upwards almost vertically at $+85.1^\circ$.
 • $\mathbf{I}_4$: Length 40.1, pointing almost horizontally leftwards into the second quadrant at $-176.6^\circ$.
 • $\mathbf{I}_5$: Length 1.4 (shortest vector), pointing down into the fourth quadrant at $-79.7^\circ$.


(iii) Determine the real power dissipated in the network.
Real power ($P$) is purely consumed by the resistors in the network. Inductors and capacitors only consume reactive power.
Identify all resistors in the network:

1. $2\ \Omega$ resistor in the top branch carrying current $\mathbf{I}_1$.
2. $8\ \Omega$ resistor in the far-right branch carrying current $\mathbf{I}_5$.
Using the formula $P = I^2 \cdot R$ with the RMS magnitudes of the currents:
$$P_{\text{total}} = P_{2\Omega} + P_{8\Omega}$$ 
$$P_{\text{total}} = (I_{1,\text{rms}})^2 \cdot 2 + (I_{5,\text{rms}})^2 \cdot 8$$ 
$$P_{\text{total}} = (10.31)^2 \cdot 2 + (1.40)^2 \cdot 8$$ 
$$P_{\text{total}} = (106.30 \cdot 2) + (1.96 \cdot 8)$$ 
$$P_{\text{total}} = 212.60 + 15.68 = \mathbf{228.28\text{ W}}$$ 

